Платёжная система выглядит монолитом только из бизнес-API. Изнутри
она разделена на компоненты, у~каждого из которых собственный
жизненный цикл, собственное состояние и собственная команда
эксплуатации. Главы не предлагают рецепта \enquote{как собрать
платёжную платформу с~нуля}: их задача~--- показать, какой компонент
за что отвечает, какие интерфейсы между ними сложились на практике
и~почему ответственность распределена именно так.

Платёжный шлюз превращает HTTPS-запрос продавца в~авторизационное
сообщение ISO~8583. Запрос проходит через приём API, оценку риска,
маршрутизацию между эквайерами, изолированное хранилище карточных
данных и~модуль авторизации; каждый следующий компонент получает
более узкое представление транзакции, чем предыдущий, и~из этого
порядка следуют требования к~PCI scope и~ограничения на~обработку
отказов.

Внутренний реестр обязательств отвечает за состояние денег внутри
продукта: что уже заработано, что зарезервировано, что можно
выплатить продавцу. Реестр опирается на двойную запись и~событийную
модель (авторизация, подтверждение, возврат, спор); в~маркетплейсе
один внешний платёж порождает несколько проводок с~разными сроками
признания и~разными правилами отмены, и~выразить это статусом одной
транзакции нельзя.

Мультивалютность~--- вопрос о~том, кто фиксирует курс: эмитент,
эквайер с~DCC, платформа или внешний FX-провайдер. От ответа
зависит, кто несёт валютный риск и~как устроена сверка между
валютами. Для российских платежей с~2022~года к~этому добавились
SWIFT-ограничения, перестройка корреспондентских схем через банки
дружественных стран и~рост доли юаня.

Процессинговый центр стоит между шлюзом и~карточной сетью:
коммутатор авторизации с~таблицей BIN, пул HSM, клиринговый
и~расчётный модули. Архитектурное решение проверяется тем, что
произойдёт с~задержкой и~доступностью при отказе внешнего канала
и~включении \textit{stand-in}, и~тем, кто контролирует ключевой
материал.

У~одной компании платёжный шлюз и~внутренний реестр находятся в~одном
сервисе, у~другой~--- в~нескольких системах разных команд. Граница
между шлюзом, реестром и процессингом сдвигается с~объёмом платежей,
профилем рисков и~организационным устройством, но~сам список задач
не меняется: их кто-то всё равно решает.

После части~V книгу замыкает эпилог о~том, как удерживать актуальное
знание в~индустрии, где документы и~правила быстро устаревают.
